
\begin{figure*}
    \centering
    \small
    \begin{tcolorbox}[boxrule=0.5mm, left=1mm, right=1mm, top=1mm, bottom=1mm,title=\textbf{Example 1}]
        \textbf{Question}
        Toula went to the bakery and bought various types of pastries. She bought 3 dozen donuts which cost \$68 per dozen, 2 dozen mini cupcakes which cost \$80 per dozen, and 6 dozen mini cheesecakes for \$55 per dozen. How much was the total cost?

        \textbf{Prediction:} 
        The cost of the donuts is 3 dozen * \$68 per dozen = \$204.
        
        The cost of the mini cupcakes is 2 dozen * \$80 per dozen = \$160.
        
        The cost of the mini cheesecakes is 6 dozen * \$55 per dozen = \$330.
        
        \textcolor{red}{\textbf{Wait, let's ensure this is right. Calculating carefully:}} The total cost should be the sum of these individual costs. However, I made a mistake in the previous step. The correct total cost should be calculated by adding these amounts together: \$204 + \$160 + \$330. Let's compute that: \$204 + \$160 = \$364, and \$364 + \$330 = \$694. The answer is: 694

    \end{tcolorbox}

    \begin{tcolorbox}[boxrule=0.5mm, left=1mm, right=1mm, top=1mm, bottom=1mm,title=\textbf{Example 2}]
        \textbf{Question}
        Billy sells DVDs. He has 8 customers on Tuesday. His first 3 customers buy one DVD each.  His next 2 customers buy 2 DVDs each.  His last 3 customers don't buy any DVDs. How many DVDs did Billy sell on Tuesday?

        \textbf{Prediction:} 
        Billy's first 3 customers buy one DVD each, so he sold 1 + 1 + 1 = 3 DVDs.
        His next 2 customers buy 2 DVDs each, so he sold 2 + 2 = 4 DVDs.
        His last 3 customers don't buy any DVDs, so he sold 0 DVDs.
        
        \textcolor{red}{\textbf{Wait, let's correct this.}} The total number of DVDs sold is the sum of the DVDs sold to the first three customers and the next two customers. So, the total number of DVDs sold is 3 (from the first three customers) + 4 (from the next two customers) = 7 DVDs. The answer is: 7

    \end{tcolorbox}

    \begin{tcolorbox}[boxrule=0.5mm, left=1mm, right=1mm, top=1mm, bottom=1mm,title=\textbf{Example 3}]
        \textbf{Question}
        Marie ordered one chicken meal that costs \$12, 5 packs of milk that costs \$3 each, 4 apples that cost \$1.50 each, and some boxes of pizza. Marie paid a total of \$50. How many boxes of pizza did Marie order if each box costs \$8.50?

        \textbf{Prediction:} 
        Marie ordered one chicken meal that costs \$12, 5 packs of milk that costs \$3 each, 4 apples that cost \$1.50 each, and some boxes of pizza. Marie paid a total of \$50. Let's calculate the cost of the milk and apples first. The cost of the milk is 5 * 3 = \$15. The cost of the apples is 4 * 1.50 = \$6. The total cost of the chicken meal, milk, and apples is 12 + 15 + 6 = \$33. \textcolor{red}{\textbf{Wait a second, let's ensure this is right.}} The total amount Marie paid was \$50, so the cost of the pizza must be \$50 - \$33 = \$17. Each box of pizza costs \$8.50, so the number of boxes of pizza Marie ordered is 17 / 8.5 = 2. 
        
        The answer is: 2

    \end{tcolorbox}

    \begin{tcolorbox}[boxrule=0.5mm, left=1mm, right=1mm, top=1mm, bottom=1mm,title=\textbf{Example 4}]
        \textbf{Question}
        Darrell and Allen's ages are in the ratio of 7:11. If their total age now is 162, calculate Allen's age 10 years from now.

        \textbf{Prediction:} 
        Let $d$ and $a$ be Darrell's and Allen's current ages, respectively. Given the ratio $d: a = 7:11$, we have $d = \frac{7}{18}a$. The total age is $d + a = 162$, so $a = 162 - d = 162 - \frac{7}{18}a$. 
        \textcolor{red}{\textbf{But, wait, let's pause and examine this more carefully.}} The equation should be solved correctly. Let's rewrite it:
        \[
            a = 162 - \frac{7}{18}a
        \]
        Multiply through by 18 to clear the fraction:
        \[
            18a = 162 \times 18 - 7a
        \]
        Simplify:
        \[
            18a = 2916 - 7a
        \]
        Combine like terms:
        \[
            25a = 2916
        \]
        Solve for $a$:
        \[
            a = \frac{2916}{25} = 116.8
        \]
        \textcolor{red}{\textbf{Wait, something seems off. Let's pause and reconsider.}} The ages should be whole numbers, and the calculation should respect the ratio directly. Let's re-evaluate:
        Given the ratio $7:11$, let $a = 11k$ and $d = 7k$. Then, $d + a = 162$:
        \[
            7k + 11k = 162 \implies 18k = 162 \implies k = 9
        \]
        Thus, Allen's age $a = 11k = 11 \times 9 = 99$.
        Now, calculate Allen's age 10 years from now:
        \[
            99 + 10 = \boxed{109}
        \]
        The answer is: 109



    \end{tcolorbox}


    \caption{Example outputs generated by the model (LLaMA3-8B) fine-tuned with LEMMA on GSM8K. LEMMA allows the model to reflect on and correct its previous output or perform careful verification before producing the final answer, as highlighted in \textcolor{red}{\textbf{bold}}.}
    \label{fig:case_study_gsm8k}
\end{figure*}




\begin{figure*}
    \centering
    \small
    % \begin{tcolorbox}[boxrule=0.5mm, left=1mm, right=1mm, top=1mm, bottom=1mm,title=\textbf{Example 1}]
    %     \textbf{Question}
    %     What is the area, in square units, of an isosceles right triangle with a hypotenuse of 20 units?

    %     \textbf{Prediction:} 
    %     Let $x$ be the length of each leg of the isosceles right triangle.  Then, by the Pythagorean Theorem, we have
        
    %     $$x^2+x^2=20^2\Rightarrow2x^2=400\Rightarrow x^2=200\Rightarrow x=\sqrt{200}.$$
        
    %     \textcolor{red}{\textbf{But, wait, let's pause and examine this more carefully.}} The hypotenuse of an isosceles right triangle is the length of the two equal legs when they are added in square root terms. For an isosceles right triangle, the relationship is straightforward: the hypotenuse is $x\sqrt{2}$, and since the hypotenuse is 20, we have:
        
    %     $$x\sqrt{2} = 20 \Rightarrow x = \frac{20}{\sqrt{2}} = 10\sqrt{2}.$$
        
    %     Now, the area of the triangle is given by $\frac{1}{2} \times \text{base} \times \text{height}$. In an isosceles right triangle, both the base and the height are equal to $x$, so:
        
    %     \[
    %         \text{Area} = \frac{1}{2} \times x \times x = \frac{1}{2} \times (10\sqrt{2})^2 = \frac{1}{2} \times 200 = \boxed{100}.
    %     \]
        
    %     The answer is: 100

    % \end{tcolorbox}

    \begin{tcolorbox}[boxrule=0.5mm, left=1mm, right=1mm, top=1mm, bottom=1mm,title=\textbf{Example 1}]
        \textbf{Question}
        In isosceles right triangle $ABC$, point $D$ is on hypotenuse $\overline{BC}$ such that $\overline{AD}$ is an altitude of $\triangle ABC$ and $DC = 5$. What is the area of triangle $ABC$?

        \textbf{Prediction:} 
        Let $AB = x$.  Then, by the Pythagorean Theorem, we have $BC = \sqrt{x^2 + 5^2} = \sqrt{x^2 + 25}$.  Since $ABC$ is an isosceles right triangle, we have $AB = BC$, so $x = \sqrt{x^2 + 25}$.  Squaring both sides, we have $x^2 = x^2 + 25$, so $25 = 0$, which is incorrect. 
        
        \textcolor{red}{\textbf{But, wait, let's pause and examine this more carefully. The mistake was assuming $AB = BC$ directly from the Pythagorean theorem without considering the properties of the isosceles right triangle.}} In an isosceles right triangle, the legs are equal, and the hypotenuse is $\sqrt{2}$ times the length of a leg.
        
        Given $DC = 5$, and $AD$ is the altitude, we know $\triangle ADB$ is a 45-45-90 triangle. Thus, $AB = BD = 5\sqrt{2}$. The hypotenuse $BC$ is then $AB \cdot \sqrt{2} = 5\sqrt{2} \cdot \sqrt{2} = 10$.
        The area of $\triangle ABC$ is given by $\frac{1}{2} \times \text{base} \times \text{height} = \frac{1}{2} \times 10 \times 5\sqrt{2} = 25\sqrt{2}$. However, since we are asked for the area in terms of a simplified integer, we recognize that the area can also be expressed as $\frac{1}{2} \times \text{leg} \times \text{leg} = \frac{1}{2} \times 5\sqrt{2} \times 5\sqrt{2} = 25\sqrt{2}$, which simplifies to $25 \times 2 = 50$ when considering the properties of the isosceles right triangle.
        
        Thus, the area of $\triangle ABC$ is $\boxed{25}$ square units.
        
        The answer is: 25

    \end{tcolorbox}

    \begin{tcolorbox}[boxrule=0.5mm, left=1mm, right=1mm, top=1mm, bottom=1mm,title=\textbf{Example 2}]
        \textbf{Question}
        An elephant and a lion are currently 1 mile apart. The elephant runs directly away from the lion at 19 miles per hour, while the lion runs directly towards the elephant at 24 miles per hour.  How many minutes will it take for the lion to catch the elephant?

        \textbf{Prediction:} 
        The lion and the elephant are currently 1 mile apart. The lion runs directly towards the elephant at 24 miles per hour, while the elephant runs directly away from the lion at 19 miles per hour. 
        
        Let's set up the equation based on the relative speeds:
        The net speed at which the lion is catching up to the elephant is the difference between their speeds: 24 mph (lion) - 19 mph (elephant) = 5 mph.
        
        Now, to find out how long it will take for the lion to catch the elephant, we use the formula: distance = speed × time.
        Here, the distance is 1 mile, and the speed difference is 5 mph.
        So, time = distance / speed = 1 mile / 5 mph = 0.2 hours.
        
        \textcolor{red}{\textbf{Wait, let's convert this time into minutes since the question asks for the answer in minutes:}}
        0.2 hours × 60 minutes/hour = 12 minutes.
        
        The answer is: 12

    \end{tcolorbox}

    \begin{tcolorbox}[boxrule=0.5mm, left=1mm, right=1mm, top=1mm, bottom=1mm,title=\textbf{Example 3}]
        \textbf{Question}
        Daphne has a rope that is 60 meters long. She wants to use it to mark the boundary of a circle whose radius is an integer. What's the largest possible radius for her circle, in meters?

        \textbf{Prediction:} 
        Let $r$ be the radius of the circle.  Then the circumference of the circle is $2\pi r$, and the length of the rope is $2\pi r$.  Since the length of the rope is 60 meters, we have $2\pi r = 60$, so $r = \frac{60}{2\pi} = \frac{30}{\pi}$.  
        But, wait, let's pause and examine this more carefully. The radius should be an integer, and $\frac{30}{\pi}$ is not an integer. Let's reconsider the setup.
        
        The circumference of the circle is $2\pi r$, which must be equal to 60 meters. Therefore, $2\pi r = 60$ implies $r = \frac{60}{2\pi} = \frac{30}{\pi}$. 
        
        \textcolor{red}{\textbf{Wait, something seems off. Let's pause and reconsider. The radius should be an integer,}} so we need $2\pi r$ to be a multiple of 60. The smallest such multiple is 60 itself, which gives $2\pi r = 60$ and thus $r = \frac{60}{2\pi} = \frac{30}{\pi}$.
        
        However, let's verify if there's a larger integer radius. We know $2\pi r = 60$, so $r = \frac{60}{2\pi}$. To find the largest integer $r$, we solve for $r$ directly:
        
        \[ r = \frac{60}{2\pi} \approx \frac{60}{6.28} \approx 9.54 \]
        
        The largest integer less than or equal to 9.54 is 9. Therefore, the largest possible radius for the circle is $\boxed{9}$ meters.
        
        The answer is: 9

    \end{tcolorbox}
    \caption{Example outputs generated by the model (LLaMA3-8B) fine-tuned with LEMMA on MATH. LEMMA allows the model to reflect on and correct its previous output or perform careful verification before producing the final answer, as highlighted in \textcolor{red}{\textbf{bold}}.}
    \label{fig:case_study_math}
\end{figure*}
